% ============================================================
%  CHAPTER 1 — INTRODUCTION
% ============================================================
\chapter{Introduction}

\section{Motivation}

The rapid shift toward remote learning, online education, and digital workplaces
has created an urgent need for tools that can objectively measure learner and
employee engagement. In a traditional classroom or meeting room, an instructor or
manager can intuitively observe whether participants are attentive, confused,
frustrated, or simply disengaged. However, in virtual environments — or even
large physical classrooms — such qualitative assessment becomes impractical at
scale.

Research consistently shows that emotional state is one of the strongest
predictors of learning outcomes and work productivity. A student experiencing
confusion or frustration is likely to disengage if not promptly supported,
while a student expressing curiosity or happiness is more likely to retain
information. Despite this, most e-learning and remote work platforms today
provide no mechanism for real-time affective feedback.

Advances in deep learning, particularly in facial action recognition and transfer
learning, now make it computationally feasible to detect human emotions in real
time using only a standard webcam. MindTrace leverages these advances to build a
non-invasive, automated engagement tracking system that works entirely through the
user's existing camera, requiring no wearable sensors, specialised hardware, or
manual input.

\section{Problem Definition and Objectives}

\subsection{Problem Definition}

To design and implement a real-time emotion recognition and engagement tracking
system using deep learning that analyses a user's facial expressions through a
webcam, computes a continuous focus/engagement score, persists session-level
analytics in a cloud database, and presents insights through an interactive web
dashboard — all without requiring any specialised hardware beyond a standard
computer camera.

\subsection{Objectives}

\begin{itemize}
    \item To detect and classify the emotional state of a user in real time from
          webcam frames into seven categories: \textit{Angry, Disgust, Fear,
          Happy, Neutral, Sad,} and \textit{Surprise}.

    \item To estimate head pose (pitch, yaw, roll) and detect attentional
          distraction when a user looks away from the screen.

    \item To compute a composite engagement/focus score (0--100) by combining
          emotional state, head pose, and blink rate signals.

    \item To persist all session data, emotion logs, and engagement metrics in a
          MongoDB database for historical analysis and reporting.

    \item To provide a React-based analytics dashboard offering live emotion
          overlays, session timelines, emotion distribution charts, and an admin
          panel for multi-user monitoring.

    \item To secure the system with JWT-based user authentication and bcrypt
          password hashing, supporting role-based access control for students
          and administrators.

    \item To deploy the backend on a cloud platform (Hugging Face Spaces) and
          the frontend on Vercel, enabling fully remote, hardware-independent
          operation.
\end{itemize}

\section{Project Scope}

MindTrace is designed for deployment in the following contexts:

\begin{itemize}
    \item \textbf{Educational Institutions} — Monitoring student engagement
          during online lectures, self-paced video courses, or e-assessment
          sessions to identify students at risk of disengagement.

    \item \textbf{Corporate Training} — Tracking employee attention and
          emotional response during virtual onboarding, compliance training,
          or skill-development workshops.

    \item \textbf{Remote Work Environments} — Providing managers with
          aggregated, anonymised team engagement metrics during video
          conferences or collaborative sessions.

    \item \textbf{Research} — Serving as a data collection and annotation
          platform for affective computing researchers studying the relationship
          between emotion and cognitive load.
\end{itemize}

The system operates entirely client-side for frame capture and relies on a
RESTful FastAPI backend for model inference, keeping user data private and
processing load distributed. The modular architecture allows individual
components — the emotion model, engagement algorithm, or frontend dashboard —
to be upgraded independently as the project evolves.
